\begin{afterword}
\summaryhead

The last post of the book looks back on what the author has written about rationality and says that more of the art is missing than present: defeating akrasia, coordinating groups, training, teaching and verification among it. The author guesses that there is a barrier to getting started. People who begin to think about thinking go wrong by default, as David Stove said of the great philosophers such as Hegel, and end up with something like Freudian psychoanalysis or a new religion. The hope is that the part of the art already written will get readers past that barrier and give them a ``saving throw.''

The rest of the work is left to readers, with two pieces of advice, both hedged: draw on many sources, not one author, and develop the art while trying to do something that matters. The post ends by asking readers to go out, achieve something beyond the armchair, create new art, and report back.

\responsehead

The post is an exhortation that closes a book, and it should be judged as one. It does not argue for a thesis; it takes stock, states a hope and hands on a task. Judged that way, it has little to fault.

Its account of what the art lacks is candid. It says that more is absent than present, that the author has worked more on belief than on action, and that training, teaching and verification are still missing. These points gather what earlier posts in the sequence argued, and the notes point back to them rather than repeat them.

Its one argument, that people who set out to develop an art of thinking go wrong by default, is labelled a guess and credited to David Stove. The attribution holds. Stove's essay calls the ``great thinkers'' ``proper objects, indeed, of pity,'' and says that people who attempt ``any depth or generality of thought'' go mad ``almost infallibly.'' The evidence behind the guess is Stove's examples and the post's own mention of Freud, not a survey, and the post does not claim more.

Its second piece of advice is more careful than the book's earlier version of it. ``Something to Protect'' said without a hedge that no one masters the Way until more than their life is at stake. Here the author wonders whether this is ``overgeneralizing from my own history'' and names the alternative, that the next advances may come from direct study of rationality.

For information, not as a criticism: the task the post hands on was taken up. The Center for Applied Rationality was founded in 2012 to teach rationality in workshops. What its own evaluation found is reported in our notes on ``3 Levels of Rationality Verification.''

In short: a closing exhortation that is candid about what the art lacks, credits its guess about why thinkers go wrong accurately to David Stove, and hedges its lesson about purpose more carefully than ``Something to Protect'' did.
\end{afterword}
